% !TEX program = xelatex
\documentclass[11pt,a4paper]{report}

\usepackage[a4paper,margin=27mm]{geometry}
\usepackage{fontspec}
\usepackage{microtype}
\usepackage{setspace}
\usepackage{booktabs}
\usepackage{longtable}
\usepackage{tabularx}
\usepackage{array}
\usepackage{enumitem}
\usepackage{hyperref}
\usepackage{xcolor}
\usepackage{listings}
\usepackage{titlesec}

\setmainfont{Latin Modern Roman}
\setsansfont{Latin Modern Sans}
\setmonofont{Latin Modern Mono}
\onehalfspacing
\setlist[itemize]{leftmargin=*,itemsep=2pt}
\setlist[enumerate]{leftmargin=*,itemsep=2pt}
\hypersetup{colorlinks=true,linkcolor=blue!45!black,urlcolor=blue!55!black,citecolor=blue!45!black}
\renewcommand{\arraystretch}{1.18}
\setcounter{tocdepth}{2}

\title{\textbf{GenTL Swift Consumer Architecture}\\
\large Engineering Research Record, Verification Model, and Test Strategy\\[1em]
\normalsize Internal research and implementation document}
\author{AMSCOPE\_DRIVER}
\date{20 November 2026}

\begin{document}
\maketitle

\begin{abstract}
This document records the engineering model, architectural decisions, assumptions, verification strategy, known limitations, and acceptance criteria for the Swift GenTL Consumer implementation in AMSCOPE\_DRIVER. It is an internal research and development record, not a publication, certification claim, or assertion of completed conformance.

The central design is one application-facing Swift protocol, \texttt{GenTLAPI}, implemented by two independent backends: a native GenTL Consumer that dynamically loads manufacturer-supplied CTI Producer libraries, and an Amcam SDK adapter that maps the same contract to the Amcam SDK where the SDK provides equivalent capabilities. The native path must cross the actual GenTL C ABI; the Swift protocol is not a Producer and does not replace the CTI.

A physical GenTL-capable camera is not available for current development. This does not prevent rigorous testing of the protocol contract, C ABI declarations, dynamic loading, symbol resolution, dispatch, error propagation, router lifecycle, and a substantial simulated acquisition path. A deterministic native Mock GenTL Producer is therefore specified as a first-class test instrument. Hardware-dependent claims remain explicitly outside the evidence obtainable from the mock.
\end{abstract}

\tableofcontents

\chapter{Purpose, Scope, and Evidence Discipline}

\section{Purpose}
The purpose is to establish a stable, testable boundary between application code and camera backends while retaining actual GenTL interoperability. The work must avoid parallel camera abstractions, duplicated ABI declarations, fabricated backend capabilities, and application rewrites caused by backend implementation details.

This record serves four functions:
\begin{enumerate}
  \item preserve the decisions and their technical reasons;
  \item define the evidence required to establish correctness;
  \item make testing possible without physical GenTL hardware;
  \item distinguish source-level implementation from compiler, runtime, and hardware verification.
\end{enumerate}

\section{Scope}
In scope:
\begin{itemize}
  \item the pinned GenTL 1.6 C header and its imported ABI types;
  \item the 64 GenTL function operations currently represented by \texttt{GenTLAPI};
  \item native CTI loading and C function-pointer dispatch;
  \item the Amcam SDK-backed implementation of the same Swift protocol;
  \item backend selection, discovery, ownership, and lifecycle in \texttt{GenTLRouter};
  \item deterministic mock Producer design and automated test strategy;
  \item simulated discovery, device opening, stream, buffer, event, and acquisition behavior.
\end{itemize}

Out of scope for a claim of completion:
\begin{itemize}
  \item GenTL certification;
  \item universal compatibility with every vendor CTI;
  \item proof of real transport reliability, sensor timing, trigger behavior, or device firmware behavior without hardware;
  \item unverified GenApi XML, register maps, chunk data, multipart payloads, or vendor-specific features;
  \item replacing the existing application-facing Swift protocol or creating a second camera abstraction.
\end{itemize}

\section{Evidence labels}
All project reports should use the following evidence distinctions:
\begin{description}[style=nextline]
  \item[Source inspected] Code or declarations were reviewed. This does not establish compilation or runtime behavior.
  \item[Compiled] The relevant package or target compiled with a recorded toolchain and command.
  \item[Unit tested] A test executed and passed; the exact test target and environment are recorded.
  \item[Integration tested] Multiple real components interacted in the test environment, for example Swift native loader to a compiled mock CTI.
  \item[Hardware tested] A named physical device, transport, Producer version, operating system, and test conditions were used.
  \item[Externally conformant] Only use when a defined external conformance procedure and its evidence support the claim.
\end{description}
A source-level signature comparison must never be described as ABI runtime verification. A passing mock test must never be described as hardware validation.

\chapter{System Model and Architectural Boundaries}

\section{System roles}
The architecture has one shared contract and two interchangeable implementations.

\subsection{\texttt{GenTLABI}}
A standalone Swift package containing the single authoritative, pinned GenTL C header and ABI layout checks. It is the source of imported C types, constants, enumerations, structures, opaque handles, and native function-pointer typedefs. The header must not be copied into multiple packages.

\subsection{\texttt{GenTLAPI}}
The only application-facing GenTL contract. It is a Swift protocol that mirrors the GenTL C function surface using imported ABI types, raw pointers, explicit output parameters, and \texttt{GC\_ERROR} return values.

The protocol is not:
\begin{itemize}
  \item a C-exporting GenTL Producer;
  \item a replacement for a vendor CTI;
  \item a separately invented camera abstraction;
  \item permission to change GenTL error, handle, or output-parameter semantics.
\end{itemize}

The currently inspected source declares 64 operations. A source-level comparison against the pinned header reported matching names, parameter types, and return types. Compiler and runtime verification remain separate required gates.

\subsection{\texttt{GenTLNative}}
A native GenTL Consumer backend. It dynamically loads a manufacturer-provided \texttt{.cti} library and resolves its exported GenTL functions. It calls the actual native function pointers using typedefs imported from \texttt{GenTLABI}. It must not simulate Producer behavior in Swift or manually maintain a divergent function signature table.

\subsection{\texttt{AmcamGenTL}}
An independent implementation of \texttt{GenTLAPI} that maps supported operations to the Amcam SDK. It does not produce a CTI library and must not impersonate a native GenTL Producer. Unsupported capabilities must return honest GenTL errors. The backend must not invent GenApi XML, register maps, chunk data, multipart payloads, or other unsupported functionality.

\subsection{\texttt{GenTLRouter}}
Registers and selects backend implementations, owns backend contexts, aggregates discovery, and binds candidates and sessions to the backend that created them. Application code uses the shared contract and must not depend on vendor SDK implementation types.

\section{Call paths}
The native path is:

\begin{verbatim}
Application
  -> GenTLRouter
      -> GenTLAPI (Swift protocol contract)
          -> GenTLNativeImplementation
              -> imported C function-pointer typedef
                  -> dlopen/dlsym-resolved CTI export
                      -> vendor GenTL Producer
                          -> transport / device
\end{verbatim}

The Amcam path is:

\begin{verbatim}
Application
  -> GenTLRouter
      -> GenTLAPI (same Swift protocol contract)
          -> AmcamGenTL
              -> Amcam SDK
                  -> Amcam camera / driver
\end{verbatim}

The contract is shared; backend execution mechanisms are not. Native interop is genuine C ABI forwarding. Amcam is a direct SDK mapping.

\section{Package dependency constraints}
The intended dependency graph is:

\begin{verbatim}
                 GenTLABI
                    |
                 GenTLAPI
                 /      \
        GenTLNative    AmcamGenTL
                 \      /
                 GenTLRouter
                      |
                 Application
\end{verbatim}

\texttt{GenTLNative} and \texttt{AmcamGenTL} both depend on the same ABI and API packages. The router composes the implementations. There must be one authoritative header, one application-facing contract, and no application dependency on internal migration SPI.

\chapter{Architectural Decisions and Rationale}

\begin{longtable}{p{0.07\textwidth}p{0.27\textwidth}p{0.29\textwidth}p{0.29\textwidth}}
\toprule
ID & Decision & Reason / invariant & Consequence and control \\
\midrule
D-01 & Preserve one \texttt{GenTLAPI} Swift protocol. & The application must not be rewritten when selecting a backend; both backends need one observable contract. & Backend-specific implementation remains behind protocol conformance. Avoid introducing another Camera API. \\
D-02 & Keep native CTI interop real. & GenTL is a C ABI; vendor CTIs provide Producer implementations. & Use \texttt{dlopen}/\texttt{dlsym} and imported C function-pointer typedefs. A Swift simulation cannot substitute for this test. \\
D-03 & Separate \texttt{GenTLABI}. & ABI declarations are shared by API, native, and Amcam packages. & One header location prevents drift and circular ownership. \\
D-04 & Preserve C-style results at the raw boundary. & GenTL functions return \texttt{GC\_ERROR} and use output pointers. & Do not replace the contract with throws, result wrappers, or transformed handles. Convenience layers must be separate and thin. \\
D-05 & Keep Amcam as a direct SDK backend. & Amcam does not supply a CTI in this architecture. & No synthetic CTI and no claim that it is a binary Producer. Unsupported features remain explicitly unsupported. \\
D-06 & Resolve later-revision exports according to standard revision. & A GenTL 1.6 header does not imply every Producer implements every 1.6 function. & Require the complete GenTL 1.0 base; report absent later functions as not implemented where applicable. \\
D-07 & Build a compiled native Mock Producer. & No physical CTI-capable camera is available, but loader and ABI paths still require real integration tests. & Test the actual dynamic loader and C dispatch, not only Swift mocks. Mock evidence is not hardware evidence. \\
D-08 & Use contract tests for both implementations, with capability-aware expectations. & Shared API semantics should not drift; hardware/SDK capability differences are real. & Common tests cover invariant semantics; backend-specific tests cover translation and unsupported behavior. \\
D-09 & Fault injection is a designed mock capability. & Cleanup and error paths are difficult to reproduce with hardware and timing-dependent failures. & Deterministic failure modes allow lifecycle, rollback, and recovery assertions. \\
D-10 & Keep transitional \texttt{BackendConvenience} SPI internal during migration. & Existing backend lifecycle code still depends on typed handles and throwing helpers. & Do not expose it to applications. Remove only after dependencies are migrated and tests cover equivalent behavior. \\
D-11 & Separate simulated and hardware acceptance. & A mock cannot establish physical transport, firmware, or timing correctness. & Report mock conformance and hardware validation as different milestones. \\
\bottomrule
\end{longtable}

\section{Function revision policy}
The pinned header and applicable GenTL revision inventory determine symbol classification. The current implementation plan records the GenTL 1.0 base set as required, with \texttt{EventGetDataInfo} included. Later-revision additions are resolved according to their introduction:
\begin{itemize}
  \item GenTL 1.1: port URL enumeration and stacked port register operations;
  \item GenTL 1.3: buffer chunk data;
  \item GenTL 1.4: parent-module queries;
  \item GenTL 1.5: buffer-part enumeration and information;
  \item GenTL 1.6: composite buffers, stacked buffer information, flows, and buffer segments.
\end{itemize}
The precise export inventory must be mechanically checked against the pinned header and standard revision before the mock fixture encodes it. No symbol should be marked optional merely because it is inconvenient to implement.

\chapter{Testing Research: Why Hardware Is Not a Prerequisite}

\section{Verification decomposition}
The test strategy separates concerns that are often incorrectly collapsed into a single camera test.

\begin{longtable}{p{0.25\textwidth}p{0.34\textwidth}p{0.31\textwidth}}
\toprule
Layer & Primary evidence & Requires physical camera? \\
\midrule
Swift protocol surface & Conformance, argument/output behavior, default unsupported results & No \\
Imported ABI declarations & Header-derived types, layout checks, compile-time type checking & No \\
Native dynamic loader & Library loading, symbol lookup, required/optional symbol policy & No \\
Native dispatch & Correct C function called with exact arguments and output behavior & No \\
Router & Registration, discovery, ownership, lifecycle, rollback, concurrency & No \\
Simulated stream & Buffer queues, events, frame metadata, deterministic payloads & No \\
Amcam translation & SDK seam behavior, mapping, unsupported capability handling & No for injected tests; yes for actual SDK/hardware validation \\
Vendor CTI compatibility & Vendor-specific behavior and interpretation & A vendor CTI is required; physical device may not be required for every Producer function \\
Physical transport & Actual link, sensor, trigger, timing, firmware, and image delivery & Yes \\
\bottomrule
\end{longtable}

A CTI library and a physical camera are distinct dependencies. A vendor CTI can often be loaded and queried without a connected camera. A mock CTI can be built without either. Tests must state which dependency they exercise.

\section{Why a Swift-only mock is insufficient}
A Swift fake conforming to \texttt{GenTLAPI} is useful for router unit tests, but it bypasses:
\begin{itemize}
  \item the dynamic loader;
  \item symbol names and export visibility;
  \item C calling convention and function-pointer typedefs;
  \item ABI-level pointer and structure interactions;
  \item native library lifetime and unload behavior.
\end{itemize}
Therefore it cannot be the only test double for \texttt{GenTLNative}. A compiled C shared library exporting GenTL functions is required for the native integration boundary.

\section{Mock Producer design}
The Mock Producer is a deterministic C shared library compiled as a dynamic library and loaded by the production \texttt{GenTLNative} loader. On macOS it may be emitted as a \texttt{.dylib}; the loader-facing path/name must match the package's actual CTI acceptance rules. Naming a library \texttt{.cti} does not itself make it conformant: it must export the correct ABI and behave according to the tested GenTL subset.

The fixture must:
\begin{enumerate}
  \item include the same pinned GenTL header through the shared ABI target or a controlled C fixture include path;
  \item export the required GenTL 1.0 base symbols with exact prototypes and calling conventions;
  \item expose a deterministic system, interface, device, and stream topology;
  \item implement enough standard behavior for discovery, device open/close, stream setup, buffer queueing, acquisition, and event delivery;
  \item expose test-only control functions for fault injection and call recording, isolated from the GenTL export surface;
  \item offer build variants with selected missing required or optional symbols;
  \item provide deterministic reset and cleanup behavior between test cases.
\end{enumerate}

Test-only control exports must use a separate prefix (for example \texttt{MockCTI\_*}) and must not be mistaken for GenTL API functions.

\section{Deterministic simulated device}
The baseline simulated topology is one Producer, one System, one Interface, one Device, and one Data Stream. The fixture should support configuration for multiple devices and streams once the basic path is stable.

The initial image generator should produce a deterministic byte pattern derived from frame sequence and byte offset. This permits exact payload assertions without requiring image-processing libraries. Metadata should include stable frame IDs, dimensions, pixel format, payload size, timestamps with defined mock semantics, and buffer completeness status.

Do not pretend mock timestamps reproduce hardware clocks. They are controlled test values.

\section{Native dispatch proof}
The mock records, for each GenTL operation:
\begin{itemize}
  \item operation identifier;
  \item invocation count;
  \item scalar arguments and handle identities;
  \item input buffer bytes where the operation contract makes them readable;
  \item output values written;
  \item configured return code;
  \item optional delay or synchronization barrier for concurrency tests.
\end{itemize}

Tests call through \texttt{any GenTLAPI} backed by \texttt{GenTLNativeImplementation}, then inspect the mock's test-control interface. This establishes that a Swift protocol call reached the expected native C export with the expected arguments. Avoid reading memory beyond sizes explicitly supplied by the caller.

\chapter{Test Catalogue and Coverage Model}

\section{Protocol contract suite}
All 64 operations require an explicit test disposition. For each operation, record:
\begin{itemize}
  \item operation name and owning GenTL module;
  \item introduction revision;
  \item required/optional export classification;
  \item argument forwarding assertion;
  \item output parameter assertion;
  \item return-code assertion;
  \item null/invalid-input behavior where defined;
  \item buffer sizing and retry behavior where applicable;
  \item backend capability disposition.
\end{itemize}

A test disposition may be: implemented-and-tested, explicitly unsupported, not applicable to the backend, or blocked pending a mock capability. It must not be silently omitted.

\section{API default behavior}
The default protocol implementations are expected to return \texttt{GC\_ERR\_NOT\_IMPLEMENTED} for operations not implemented by a conformer. Tests must verify:
\begin{itemize}
  \item each default returns the exact error constant;
  \item output parameters are not falsely reported as valid data;
  \item implementing one operation does not change defaults for another;
  \item protocol extension dispatch does not accidentally bypass the concrete implementation.
\end{itemize}

\section{Native loader and ABI tests}
Required cases:
\begin{enumerate}
  \item valid mock with all required exports loads and initializes;
  \item each required export missing individually causes deterministic rejection;
  \item multiple missing required exports produce a useful diagnostic;
  \item optional later-revision export absent does not reject a valid base Producer;
  \item invoking an absent optional operation returns the documented not-implemented status;
  \item invalid file path, non-library file, incompatible architecture, and failed symbol lookup are reported safely;
  \item repeated load/close cycles do not leak library handles;
  \item failed initialization attempts perform the correct rollback;
  \item function-pointer table entries correspond to the correct symbols and imported typedefs.
\end{enumerate}

\section{Argument and output semantics}
For every function, use distinct sentinel values where possible so argument swaps are detectable. For output parameters, initialize storage to sentinel values and assert the Producer writes only the expected output. Test nullability only where the C contract permits null. Test buffer-too-small behavior with exact required-size reporting where specified. Avoid assumptions that every GenTL function uses identical string or buffer conventions.

\section{Router tests}
Router tests should cover:
\begin{itemize}
  \item duplicate and malformed backend identifiers;
  \item repeated discovery and refresh;
  \item multiple registered native Producers;
  \item mixed native and Amcam registrations;
  \item candidates remaining bound to their originating context;
  \item open session surviving a discovery refresh;
  \item close ordering and idempotency;
  \item initialization rollback;
  \item failed close/quarantine behavior;
  \item concurrent discovery and independent session operations;
  \item no use-after-close or accidental backend cross-routing.
\end{itemize}

\section{Acquisition integration tests}
The mock-backed end-to-end path should cover:
\begin{enumerate}
  \item initialize Producer and enumerate System;
  \item enumerate and open Interface;
  \item enumerate and open Device;
  \item obtain the remote Port and Data Stream identifiers;
  \item open Data Stream;
  \item allocate/announce buffers and verify ownership;
  \item queue buffers and start acquisition;
  \item receive an event and retrieve event data;
  \item verify frame metadata, payload size, deterministic payload bytes, and completeness;
  \item requeue buffers and acquire multiple frames;
  \item stop acquisition and flush queues;
  \item revoke buffers and verify returned user/private pointers;
  \item close stream, device, interface, System, and Producer in valid order.
\end{enumerate}

Tests must include invalid ordering and failure at every lifecycle boundary. The expected cleanup sequence must be explicit; a failed close must not lead to unsafe library unloading while live native handles may still reference Producer code.

\section{Fault injection matrix}
The Mock Producer should support deterministic fault injection for at least:
\begin{longtable}{p{0.34\textwidth}p{0.30\textwidth}p{0.26\textwidth}}
\toprule
Failure point & Expected assertion & Primary owner \\
\midrule
Producer init & Error preserved; no false initialized state & Native / Router \\
System open & Output handle and rollback are correct & Native / Router \\
Interface enumeration & Error and partial results handled deterministically & Router \\
Interface open & No leaked System/Interface resources & Router \\
Device open & No candidate/session state corruption & Router \\
Stream open & Device remains closable & Router / Application \\
Buffer announce/allocate & No orphaned buffer registration & Native / Application \\
Queue buffer & Ownership state remains consistent & Producer / Application \\
Start acquisition & No false running state & Producer / Application \\
Event wait & Timeout distinguished from other errors & Native / Application \\
Event payload & Buffer-too-small retry and size semantics preserved & Native / Application \\
Incomplete frame & Status and metadata remain observable & Application \\
Stop acquisition & Cleanup can continue or quarantine safely & Application / Router \\
Revoke buffer & Returned pointers and ownership are correct & Producer / Application \\
Close stream/device/interface & Failed cleanup never unloads live code & Router / Native \\
Library close & Backend cannot be reinitialized unsafely & Router \\
\bottomrule
\end{longtable}

\section{Stress and concurrency}
Stress suites should be separate from deterministic unit tests so failures remain diagnosable. Start with bounded, repeatable loops, then scale:
\begin{itemize}
  \item repeated init/terminate and load/unload;
  \item repeated discovery and open/close;
  \item multiple devices and streams;
  \item concurrent calls on independent sessions;
  \item concurrent discovery while sessions remain open;
  \item event timeout and cancellation races;
  \item buffer queue/requeue under sustained acquisition;
  \item failure injection during cleanup;
  \item sanitizer-enabled C fixture builds where supported.
\end{itemize}
Concurrency tests must not assert undocumented thread safety from the GenTL standard. They should verify the synchronization guarantees the implementation itself promises and serialize calls where the Producer contract requires it.

\chapter{Amcam Backend Verification}

\section{Contract equivalence without false capability equivalence}
Both backends implement the same protocol, but they do not necessarily support the same feature set. Tests must separate:
\begin{itemize}
  \item shared contract invariants: error model, output parameter discipline, handle validity, lifecycle behavior;
  \item Amcam mapping: a GenTL operation maps to a documented SDK operation with correct units, ranges, and ownership;
  \item unsupported behavior: unsupported GenTL features return the selected precise error and are not fabricated;
  \item hardware behavior: actual SDK calls and device behavior require an Amcam SDK environment and, for device-specific behavior, hardware.
\end{itemize}

\section{Recommended test seam}
Where the current implementation directly loads or calls the Amcam SDK, introduce a narrow internal injectable boundary around the SDK function table only if required to make deterministic unit tests possible. This seam must not become a second public camera API and must not alter the application-facing \texttt{GenTLAPI}. The production implementation must still call the real Amcam SDK.

Tests should validate handle-token lifecycle, buffer layout, event payload retry semantics, acquisition state transitions, error translation, and unsupported capability reporting. Existing source-level tests for frame layout and lifecycle should be preserved and expanded rather than duplicated.

\chapter{Risks, Limits, and Controls}

\begin{longtable}{p{0.25\textwidth}p{0.38\textwidth}p{0.27\textwidth}}
\toprule
Risk & Why it matters & Control \\
\midrule
Mock accidentally diverges from standard & Passing tests could validate a fictional Producer. & Compile fixture against the same pinned header; review each behavior against standard semantics. \\
Mock tests bypass actual loader & Native ABI bugs remain undetected. & Load the compiled shared library through production \texttt{GenTLNative}. \\
Only happy paths are tested & Resource leaks and stale handles appear under failure. & Fault injection at every lifecycle transition. \\
Protocol tests are name-only & Argument order, output pointers, and errors may still be wrong. & Per-operation sentinel forwarding and output assertions. \\
False Amcam feature parity & Application may rely on nonexistent camera functionality. & Capability-aware test matrix and honest unsupported results. \\
Mock success interpreted as hardware success & Physical transport, firmware, and timing remain untested. & Evidence labels and separate hardware acceptance. \\
Unsafe unload after failed close & Native function pointers or handles may refer to unloaded code. & Retain/quarantine provider lifetime on uncertain cleanup. \\
Test fixture overfits one device topology & Router assumptions may fail with real multi-device Producers. & Configurable topology and later multi-interface/multi-device tests. \\
\bottomrule
\end{longtable}

\chapter{Current State and Verification Record}

\section{Source-inspection baseline}
At the time this document is created, the repository source and documentation indicate:
\begin{itemize}
  \item standalone \texttt{GenTLABI} exists and is shared by the packages;
  \item \texttt{GenTLAPI} declares 64 operations using imported ABI types and C-style results;
  \item a source-level comparison against the pinned header reports matching operation names, parameter types, and return types;
  \item \texttt{GenTLNative} contains implementations and a dynamic function table based on imported C function-pointer typedefs;
  \item \texttt{AmcamGenTL} implements the common protocol and explicitly rejects unsupported capabilities;
  \item \texttt{GenTLRouter} exposes candidates and sessions bound to backend contexts;
  \item transitional \texttt{BackendConvenience} SPI remains for internal migration;
  \item a default unsupported-operation test and earlier lifecycle/frame-layout tests are documented in the implementation record.
\end{itemize}

\section{What is not established}
No claim is made here that the current branch has been built, that its tests have passed, that a compiled Mock Producer exists, or that any real Producer or physical camera has been exercised. The source-level inventory is a starting point for test implementation, not a completed verification report.

\section{Verification ledger}
Update this table only with executed evidence:
\begin{longtable}{p{0.34\textwidth}p{0.19\textwidth}p{0.37\textwidth}}
\toprule
Verification item & Status at record creation & Evidence to record \\
\midrule
Header provenance/hash and revision inventory & Source/document review required & Header path, upstream source, license, SHA-256, revision map \\
64-operation signature audit & Source inspected & Audit script/output and pinned header hash \\
Swift package compilation & Not verified & Xcode/Swift version, OS, command, result \\
ABI layout/runtime checks & Not verified & Test output and platform \\
Mock Producer compilation & Not implemented / not verified & Compiler command, artifact architecture, exported symbol list \\
Native loader integration & Not verified & Test report with mock fixture version \\
64-operation native dispatch coverage & Not verified & Per-operation test matrix and results \\
Router lifecycle/fault tests & Not verified as a complete suite & Test output and coverage report \\
Simulated acquisition & Not verified & Captured frame assertions and logs \\
Amcam SDK-injected tests & Not verified as complete & SDK seam, test output, unsupported-operation matrix \\
Physical GenTL camera & Unavailable in current environment & Device/CTI/transport/version and test report when available \\
\bottomrule
\end{longtable}

\chapter{Definition of Done}

The engineering milestone is complete only when all applicable criteria below are evidenced:
\begin{enumerate}
  \item One pinned header is the sole ABI authority and its provenance is recorded.
  \item The 64-operation Swift contract compiles against imported C declarations.
  \item ABI layout checks pass on supported build platforms.
  \item Required and optional symbols are classified by the correct GenTL revision.
  \item The production native loader loads the compiled Mock Producer.
  \item Every required native dispatch path has argument, output, and return-code assertions.
  \item Missing required symbols reject the mock; absent optional symbols follow the defined policy.
  \item Protocol default behavior is tested for all defaulted operations.
  \item Router registration, discovery, ownership, close, rollback, and quarantine tests pass.
  \item Simulated acquisition and deterministic payload verification pass.
  \item Fault injection covers lifecycle and acquisition failure boundaries.
  \item Stress tests and supported sanitizers report no known leaks or invalid memory behavior within the tested scope.
  \item Amcam mapping and unsupported behavior have explicit tests; real SDK/hardware evidence is reported separately.
  \item The sample application can discover, open, acquire from, and close the simulated device.
  \item No application-visible duplicate camera abstraction or fabricated capability remains.
  \item Hardware compatibility limitations are stated accurately.
\end{enumerate}

\chapter{Closing Engineering Position}
The absence of a physical GenTL camera is a limitation on hardware validation, not a reason to defer rigorous engineering tests. The correct response is to construct a native, ABI-accurate, deterministic Mock Producer and use it to exercise the production dynamic-loading and dispatch path. Swift fakes remain useful for isolated router tests, but they cannot replace the compiled mock at the native boundary.

The strongest defensible outcome before hardware access is: compiler-verified contract, ABI-checked declarations, native mock integration, exhaustive dispatch and lifecycle tests, deterministic simulated acquisition, and clearly bounded claims. Vendor and physical-camera validation then become an additional compatibility layer rather than the first opportunity to discover fundamental interface defects.



\chapter{Implementation Progress: First Native Mock Boundary}

\section{Fixture now present in source}
A first C fixture has been added under
\texttt{swift/Packages/GenTLNative/Tests/MockGenTLProducer/MockGenTLProducer.c}.
It includes the pinned \texttt{GenTL\_v1\_6.h} rather than redeclaring GenTL types. The fixture exports all 47 names currently listed in \texttt{GenTLProducerLibrary.requiredEntryPoints}. Forty-six functions are deterministic stubs returning \texttt{GC\_ERR\_NOT\_IMPLEMENTED}; \texttt{GCGetInfo} is deliberately implemented with observable outputs and a distinctive return code.

The fixture has a compile-time variant, \texttt{MOCK\_OMIT\_GCGETINFO}, which omits the required export while leaving the other required exports present. This permits a targeted required-symbol rejection test rather than testing only a library with no GenTL symbols at all.

\section{Native integration tests now present in source}
\texttt{GenTLNativeIntegrationTests.swift} builds the fixture at test runtime using the macOS SDK compiler and loads the resulting dynamic library through the production \texttt{GenTLNativeImplementation} initializer. The first dispatch test invokes \texttt{GCGetInfo} through an \texttt{any GenTLAPI} reference and checks:
\begin{itemize}
  \item the distinctive native return code \texttt{-4321};
  \item the datatype written by C;
  \item the deterministic payload marker written through the output pointer;
  \item the returned output size.
\end{itemize}

Additional tests verify that omission of required \texttt{GCGetInfo} rejects the CTI and that omission of optional GenTL 1.1 \texttt{GCGetNumPortURLs} does not prevent loading; invoking the absent optional operation returns \texttt{GC\_ERR\_NOT\_IMPLEMENTED} and leaves the output sentinel unchanged.

The native test target now explicitly depends on \texttt{GenTLABI} and \texttt{GenTLAPI} so tests can inspect ABI values and call the shared protocol. It reuses the package's already-declared test target.

\section{What this progress establishes and what it does not}
A source-level comparison finds 47 required entry-point names in the fixture, matching the current required-symbol inventory. This is not compiler evidence. The tests have not yet been executed; no claim is made that the C fixture compiles, the Swift tests type-check, the loader succeeds at runtime, or the dispatch assertions pass.

This initial fixture is intentionally not a simulated camera. Most exports return \texttt{GC\_ERR\_NOT\_IMPLEMENTED}. It isolates the native ABI and loader boundary first. Subsequent fixture increments must add valid System/Interface/Device discovery and then Data Stream, buffer, event, and acquisition behavior.

The first execution gate is to run \texttt{swift test --package-path swift/Packages/GenTLNative} and \texttt{swift test --package-path swift/Packages/GenTLAPI} on macOS with Xcode command-line tools. Record the exact toolchain, architecture, command, and output before proceeding.



\section{Exhaustive operation dispatch matrix added}
The mock fixture has been expanded from the required base set to all 64 GenTL function exports declared by the pinned header. Each ordinary stub returns a distinct negative result derived from its operation index. This is intentional: if a native function-table entry is accidentally associated with the wrong export, the test receives the other function's code and fails. The fixture also exports a test-only call-count query, not part of GenTL, so the test can independently establish that each native C function was actually entered.

The dispatch test invokes 58 operations through an \texttt{any GenTLAPI} value backed by \texttt{GenTLNativeImplementation} and asserts their unique per-export return codes. \texttt{GCReadPort} and \texttt{GCWritePort} are exercised separately with sentinel handles, addresses, and memory to verify argument and buffer forwarding. Producer initialization/closure and event registration/unregistration are exercised explicitly because they have lifecycle state. Event registration returns a known non-null C handle. Event unregistration internally calls \texttt{EventKill}; the test accounts for the expected second invocation of that C export.

The optional-export build variant omits all 17 later-revision exports while retaining the 47 required GenTL 1.0 base exports. The test verifies that a missing optional operation returns \texttt{GC\_ERR\_NOT\_IMPLEMENTED} and does not modify the caller's output sentinel. The required-export negative variant omits \texttt{GCGetInfo} and expects loader rejection.

A source-level audit currently reports 64 header declarations, 64 protocol requirements, 64 mock exports, 58 unique-return dispatch assertions, two port sentinel tests, and four explicit lifecycle assertions. This audit does not establish successful compilation or runtime dispatch. All native test source remains unexecuted until the macOS CI workflow or a local Xcode toolchain runs it.



\section{Argument-forwarding sentinels}
Three native operations now have deliberate argument and memory probes beyond neutral dispatch arguments:
\begin{itemize}
  \item \texttt{GCGetInfo} checks a distinctive command value and writes a datatype, payload marker, and output size.
  \item \texttt{GCReadPort} checks a non-null opaque port handle, a distinctive 64-bit address, a non-null output pointer, and the supplied capacity before writing a deterministic 32-bit value.
  \item \texttt{GCWritePort} checks the opaque handle, address, input pointer, exact four-byte payload, and supplied size.
\end{itemize}
The complete matrix therefore has 58 unique per-operation return-code assertions, two custom port argument/memory tests, and four stateful lifecycle operations. Test-only invocation counters cover all 64 protocol operations. This proves operation-to-export mapping and selected scalar/pointer forwarding paths; it is not yet exhaustive argument permutation coverage for every operation. Later test increments must add per-operation sentinel checks for info buffers, event payloads, stream buffers, and composite/stacked structures.



\section{Optional-export absence verification}
The missing-optional-export build variant now exercises all 17 later-revision Swift protocol methods. Each call must return exactly \texttt{GC\_ERR\_NOT\_IMPLEMENTED}. For operations with output parameters, the test seeds output storage with sentinel values and verifies that absent function pointers do not modify those values. This checks both loader tolerance and the raw protocol's absent-symbol behavior across all optional operations.



\section{Required-symbol rejection matrix}
The fixture supports an independent compile-time omission switch for every one of the 47 required GenTL 1.0 entry points. The native test iterates the loader's authoritative \texttt{requiredEntryPoints} inventory, builds a variant with one required export absent, and asserts that the loader rejects that exact library with a \texttt{missingRequiredSymbol} error naming the omitted export. This is stronger than a single malformed-library test because it exercises every required-symbol lookup branch. These variants are source-authored and still require macOS execution.



\section{Pinned ABI version and layout gate}
The native integration suite now asserts the compiled header version code is \texttt{0x01060000} (GenTL 1.6.0) and that \texttt{GenTLABI\_LayoutValid()} returns success. The ABI C target also contains compile-time \texttt{\_Static\_assert} checks for packed event data, port register stack entries, chunk metadata, and GenTL 1.6 stacked buffer information structures. The runtime test verifies that the linked validation entry point is callable in the test process.



\section{CI execution evidence and current infrastructure blocker}
The macOS workflow was triggered on push (run 54, commit \texttt{8512306e3c8d0cbc731fe32869bae82fb0f3278a}). All four matrix jobs were marked failed within seconds. The GitHub job records report \texttt{runner\_id: 0}, empty step arrays, and no downloadable runner logs; rerunning the failed jobs produced the same result. The available evidence therefore indicates that no runner acquired the jobs and no Swift build or test command executed. This is not evidence that the source failed to compile or that a test assertion failed.

GitHub Actions hosted-runner provisioning or repository billing/permissions must be repaired. A subsequent run included both macOS and Ubuntu jobs, and all jobs again reported runner ID zero with no steps, so the blocker is not specific to the macOS image. Alternatively, the exact package commands must be run locally on macOS with Xcode command-line tools. Until a job contains executed steps and logs, the CI conclusion must not be interpreted as a code test result.



\section{Native argument fault-injection}
The integration suite now deliberately sends invalid sentinel inputs to the three instrumented functions. An invalid \texttt{GCGetInfo} command must return its configured custom error without changing datatype, payload, or size. An invalid \texttt{GCReadPort} address must return its configured error without changing output bytes or size. An invalid \texttt{GCWritePort} payload must return its configured error while preserving the input size. These checks verify that native error values are not normalized away and that rejected operations do not manufacture output data.



\section{Invalid dynamic-library loading}
The loader unit suite now distinguishes three filesystem/load boundaries: missing path, non-regular path, and an existing regular file whose contents are not a valid Mach-O dynamic library. The third case must report \texttt{loadFailed}, not \texttt{fileNotFound} or \texttt{notRegularFile}. This validates error classification before symbol resolution begins.



The latest test workflow attempt is run \texttt{44674821941} for commit \texttt{4c108885fb95c4bf7912da5997d0e42db2d33f43}. All five jobs (four macOS package jobs and one Ubuntu ABI/API job) again terminated with runner ID zero and no executed steps. The latest source additions therefore remain uncompiled and unexecuted.


\section{CI evidence update --- 2026-11-20}
Hosted runners are now being allocated and workflow steps execute. The Ubuntu 24.04 ABI/API job in run \texttt{44676082932} (commit \texttt{430ccf9da3b098cedb79fd932874c031943317cd}) completed successfully. It built the pinned GenTL ABI package and ran the GenTLAPI contract suite: all 9 tests passed, including the test invoking all 64 protocol default operations. This verifies the Linux C ABI module build and portable Swift API contract on Ubuntu 24.04; it does not verify the native dynamic-library loader.

The CI compiler exposed two distinct test-contract defects after the typed \texttt{GC_ERROR} correction: tests compared the imported \texttt{GC_ERROR_LIST} enum against \texttt{GC_ERROR}, and passed \texttt{nil} to non-optional raw handle parameters. The test suite now uses the C ABI shim's correctly typed not-implemented status and non-null opaque handle sentinels while preserving nullability for optional output pointers. Native integration dispatch tests received the same opaque-handle correction.

The latest source commit is \texttt{8b34d21787069109027c1e8665acad01709ac3b9}. Its workflow run \texttt{44676161807} is pending/in progress; macOS package jobs have not yet produced a verified result. The native mock-producer test still compiles with \texttt{xcrun} on macOS. Linux/Windows native-loader validation and cross-platform mock compilation remain future work. No all-platform support claim is made.


\section{C ABI boundary correction --- 2026-11-20}
Further macOS compiler diagnostics exposed a structural collision between the convenience handle/status wrappers declared in \texttt{GenTLTypes.swift} and identically named imported C ABI types re-exported by \texttt{GenTLAPI}. The native implementation now keeps its high-level wrapper-based implementation in the class body and qualifies C ABI types only in its raw \texttt{GenTLAPI} conformance extension. The redundant direct \texttt{GenTLABI} import was removed from that implementation to avoid ambiguous unqualified convenience types.

The protocol and default implementations now model output handle parameters as pointers to nullable handle values (for example, \texttt{UnsafeMutablePointer<TL\_HANDLE?>?}), matching the imported C function-pointer ABI. The Amcam and native adapters were updated to the same pointer semantics. A C ABI helper, \texttt{GenTLABI\_ErrorFromCode}, explicitly converts the official \texttt{GC\_ERROR\_LIST} enum to the ABI's \texttt{GC\_ERROR} scalar type. Amcam's duplicate raw-handle helpers were removed; the remaining helper validates token identity.

The initial Linux API test pass remains valid evidence for commit \texttt{430ccf9da3b098cedb79fd932874c031943317cd}. The latest source commit is \texttt{6364f89117ea3fc76fe4cb62663913eab100f6a1}; workflow run \texttt{44676898035} completed its Ubuntu 24.04 ABI/API job successfully, including the GenTL ABI package build and all 9 GenTLAPI tests. The four macOS package jobs remain queued, so the updated Amcam/native backends and macOS integration suite are not yet verified.



\section{Amcam macOS camera path: priority and current evidence --- 2026-11-20}

The immediate product objective is reliable operation of the AmScope MU503 camera on macOS. Cross-platform native-loader work is explicitly deferred. This does not change the shared \texttt{GenTLAPI} protocol; it changes the order of engineering work. The \texttt{AmcamGenTL} implementation is the current camera path because it maps the common Swift API to the Amcam SDK directly. It is not a GenTL Producer and does not load a vendor CTI.

A deterministic native mock SDK was added as a C dynamic library fixture. The fixture is loaded through the same \texttt{dlopen}/\texttt{dlsym} code as the vendor SDK; it is not a Swift fake and does not bypass the ABI boundary. The fixture reports the Amcam SDK version expected by the implementation and models the MU503 preview modes 2560x1922, 1280x960, and 640x480. It deliberately invokes the image callback synchronously from \texttt{Amcam\_StartPullModeWithCallback}, exercising callback reentrancy before the start call returns.

The Amcam package test job completed successfully on macOS 15.8.2 arm64 with Swift 6.1.2 in workflow run \texttt{44680315441}, source commit \texttt{7e7c10f203cb0c7f620bac9a0e730567b0f8b5f3}. The Amcam job's package-test step succeeded. Its seven scheduled test cases comprise three RGB24 layout tests and four SDK integration tests:
\begin{itemize}
\item device enumeration, model/resolution metadata, camera open, three repeated start/frame/stop cycles, synchronous callback delivery, disconnect notification, and callback cleanup;
\item exposure auto/manual transitions, exposure range validation and readback, resolution selection/readback, a captured frame at 1280x960, and rejection of resolution changes while streaming;
\item data-stream input queue, finite acquisition, new-buffer event delivery, RGB24 buffer contents and metadata, and completion of the GenTL stop/close lifecycle;
\item undersized-buffer fault injection, verifying incomplete-buffer metadata and that the consumer's memory is not overwritten.
\end{itemize}

The workflow as a whole was not green: the separate \texttt{GenTLNative} and \texttt{GenTLRouter} jobs failed. Those failures are not evidence against the Amcam package job, which completed successfully. Conversely, the Amcam mock tests are not evidence that a physical MU503 has been opened or streamed on the user's Mac.

The successful run \texttt{44680315441} includes the signed 64-bit \texttt{AmcamFrameInfoV2.timestamp} ABI correction and the undersized-buffer safety test. Negative timestamp values are treated as unavailable rather than reinterpreted as very large unsigned timestamps. The job log schedules all seven test cases and the Amcam package-test step succeeds.

The physical acceptance gate remains separate:
\begin{enumerate}
\item load the actual arm64-compatible vendor \texttt{libamcam.dylib} and verify its reported version and required symbols;
\item enumerate the connected MU503 (the prior hardware profile records VID:PID \texttt{0547:4511});
\item open the camera and read back its actual current resolution and exposure controls;
\item capture and inspect frames at the available preview modes, including sequence/timestamp metadata and row stride;
\item run sustained acquisition and repeated start/stop/disconnect cycles while measuring delivered frames, camera-side loss, host-side drops, and payload goodput.
\end{enumerate}

The previous hardware profile records the user's Mac connection as USB 2.0 High Speed at 480 Mbit/s through the observed hub path. The historical Windows observation of 101 fps at 640x480 is a performance reference, not a promise for the present Mac connection. The SDK path must first measure the actual camera and transport behavior; nominal USB signaling speed and published frame rates are not substitutes for measured image-payload goodput. Do not change exposure or gain automatically to compensate for lower transport capacity.

A macOS hardware-probe executable has now been added. It uses the real AmcamGenTL backend to enumerate the camera, open an exclusive device session, acquire finite frame batches through the GenTL data-stream buffer/event path, report frame metadata and timestamps, and optionally write the last complete RGB frame as PPM. The CI workflow now builds the Apple Silicon executable and uploads it with usage instructions as an artifact. Workflow run \texttt{44680887948}, commit \texttt{805537cdedd2855c6612cb26e12f7fc3b468c96f}, completed the Amcam test, release-build, and upload steps successfully. The build log identifies the binary as a Mach-O 64-bit arm64 executable and runs its \texttt{--help} path successfully. Artifact \texttt{amcam-camera-probe-macos-arm64} (artifact ID \texttt{14604550684}) was uploaded and expires on 2026-12-04. The overall workflow remains red only because the separate GenTLNative and GenTLRouter jobs fail. The physical camera smoke test remains outstanding.

The next development work is therefore Amcam SDK correctness, camera controls, buffer ownership, event delivery, disconnect cleanup, and execution of the hardware probe on the user's Mac. Linux/Windows native-loader expansion and platform-neutral mock compilation are deferred until the AmScope camera path is reliable on macOS.



\section{Amcam hardening revision: buffer pool, ABI strings, and control readback --- 2026-11-20}

The next implementation revision remains restricted to the macOS Amcam camera path. No cross-platform loader or mock-producer refactor was started.

\subsection{SDK fixed-width string boundary}

The Amcam device enumeration structure contains fixed 64-byte identifier and display-name arrays. Passing these fields directly to Swift \texttt{String(cString:)} assumes a NUL terminator exists within the field; if a malformed SDK record omitted it, the read could continue beyond the structure field. The C ABI helper now searches for a terminator with bounded \texttt{memchr}, checks the destination capacity, and copies only the validated bytes. Swift rejects malformed or empty identifiers with an explicit enumeration error rather than silently dropping a device.

Three ABI helper tests cover valid copying, unterminated source fields, and insufficient destination capacity. The tests compile and exercise the same C helper used by the SDK loader.

\subsection{Camera control semantics}

A successful return from \texttt{Amcam\_put\_ExpoTime} is not treated as proof that the camera applied the requested value. The session now reads the value back and raises \texttt{controlReadbackMismatch} if the observed exposure differs. The mock SDK has a deterministic fault-injection switch that returns success while intentionally ignoring an exposure write; the integration test verifies the mismatch is detected and the prior readback remains observable.

The mock camera now advertises precise frame-rate support and models the SDK's minimum, maximum, and current frame-rate options. The test exercises range discovery, 0.1 FPS quantization, a successful setter/readback, an out-of-range request, an invalid increment, and the explicit unsupported result for discrete speed-level control when its optional SDK symbols are absent.

\subsection{Multi-buffer acquisition and probe integrity}

The data-stream integration test now uses two consumer-owned buffers for a three-frame finite acquisition. It verifies the expected round-robin delivery order, user-pointer payload, per-buffer metadata, requeue ownership, and completion of the asynchronous finite-acquisition stop. This exercises the buffer-pool lifecycle rather than only a single-buffer happy path.

The hardware probe now announces up to four buffers, constrained by a 512 MiB aggregate allocation ceiling. Each New Buffer event must contain the expected two-word handle/user-pointer payload, and the handle must belong to the announced pool. Before reading image memory, the probe checks the announced buffer capacity, filled size, data size, dimensions, row padding, RGB8 packed PFNC format, and pixel-format namespace. It uses a monotonic host clock, and cleanup failures are reported and produce a failing process exit instead of being silently swallowed.

A hardening baseline completed successfully in workflow run \texttt{44682176084}, source commit \texttt{0f35d0ee87e3f100a7826846c032ceab04c29880}. The macOS Amcam package job scheduled all 10 tests: three ABI helper tests, three RGB24 layout tests, and four SDK integration tests. The \texttt{Run package tests} step succeeded. The same job built the release probe, identified it as a Mach-O 64-bit arm64 executable, ran its \texttt{--help} path, and uploaded artifact \texttt{amcam-camera-probe-macos-arm64} (artifact ID \texttt{14604159604}, SHA-256 \texttt{82d0d24d1a06a341565d775d74ae544bcb7e82c71a388a4a8dd15e951bea4bcd}, expires 2026-12-04). The overall workflow remained red because the separate GenTLNative and GenTLRouter jobs failed; the Linux ABI/API, GenTLAPI, and Amcam jobs succeeded. This verifies the hardening baseline, not physical camera operation.

After that baseline, the mock was extended to test precise frame-rate range/quantization/readback and the stream test was upgraded to a two-buffer, three-frame round-robin acquisition. The probe now uses a bounded pool of up to four buffers. Latest source commit: \texttt{d5d7138ccba49c77a40c264bd4efa25e9d5eb692}; workflow run \texttt{44682507702} contains these later changes and is queued at the time of this record update. Those later changes must not be called verified until the Amcam tests, release build, and artifact upload all succeed.

The physical MU503 acceptance test remains outstanding. The artifact from run \texttt{44682176084} predates the multi-buffer and precise-frame-rate changes; use only an artifact rebuilt from the latest passing source revision for hardware acceptance.

\end{document}
